\documentclass{article}
\usepackage{mathtools} 
\usepackage{fontspec}
\usepackage[UTF8]{ctex}
\usepackage{amsthm}
\usepackage{mdframed}
\usepackage{xcolor}
\usepackage{amssymb}
\usepackage{amsmath}


% 定义新的带灰色背景的说明环境 zremark
\newmdtheoremenv[
  backgroundcolor=gray!10,
  % 边框与背景一致，边框线会消失
  linecolor=gray!10
]{zremark}{说明}

% 通用矩阵命令: \flexmatrix{矩阵名}{元素符号}{行数}{列数}
\newcommand{\flexmatrix}[4]{
  \[
  #1 = \begin{pmatrix}
    #2_{11}     & #2_{12}     & \cdots & #2_{1#4}   \\
    #2_{21}     & #2_{22}     & \cdots & #2_{2#4}   \\
    \vdots      & \vdots      & \ddots & \vdots     \\
    #2_{#31}    & #2_{#32}    & \cdots & #2_{#3#4}
  \end{pmatrix}
  \]
}

% 简化版命令（默认矩阵名为A，元素符号为a）: \quickmatrix{行数}{列数}
\newcommand{\quickmatrix}[2]{\flexmatrix{A}{a}{#1}{#2}}

\begin{document}
\title{注释 12.3}
\author{张志聪}
\maketitle

\begin{zremark}
  证明：莱布尼茨判别法的余项估计式为：
  \begin{align*}
    |R_n| \leq u_{n + 1}
  \end{align*}
\end{zremark}
证明：
\begin{itemize}
  \item $n + 1$是奇数；
        \begin{align*}
          R_n & = S - S_n                                          \\
              & = u_{n + 1} - u_{n+2} + u_{n+3} - u_{n+4} + \cdots \\
              & = \sum\limits_{k = n + 1} (-1)^{k - 1} u_k
        \end{align*}
        这里$\sum\limits_{k = n + 1} (-1)^{k - 1} u_k$也是一个
        交错级数，按照定理12.11证明类似的处理，考察它的部分和数列
        $\{T_n\}$，由区间套$\{[T_{2m}, T_{2m - 1}]\}$，且
        $T_{2m - 1}$单减，$T_{2m}$单增。于是$u_{n + 1}$就是最大的
        奇数项部分和，所以
        \begin{align*}
          u_{n + 1} \geq T_{2m - 1}
        \end{align*}
        又有
        \begin{align*}
          R_n = \lim\limits_{m \to \infty} T_{2m - 1}
        \end{align*}
        由保不等式可知
        \begin{align*}
          u_{n + 1} \geq R_n
        \end{align*}

  \item $n + 1$是偶数.
        % \begin{align*}
        %   R_n & = S - S_n                                          \\
        %       & = -u_{n + 1} + u_{n+2} - u_{n+3} + u_{n+4} + \cdots \\
        %       & = \sum\limits_{k = n + 1} (-1)^{k - 1} u_k
        % \end{align*}

        todo
\end{itemize}

\begin{zremark}
  证明：定理12.16
\end{zremark}
\textbf{证明：}
$\lim\limits_{n \to \infty} a_n = 0$，
由数列的柯西准则可知，对任意$\epsilon > 0$，存在
$N > 0$，当$n > N$时，
\begin{align*}
  |a_n - 0| & < \epsilon \\
  |a_n|     & < \epsilon
\end{align*}

$sum b_n$的部分和有界，
设部分和为$B_n$，于是存在$M > 0$，对任意$n$有
\begin{align*}
  |B_n| \leq M
\end{align*}
对任意$n, n + p$，我们有
\begin{align*}
  |B_{n + p} - B_n|
   & = \left| \sum\limits_{k = n + 1}^{n + p} b_k \right| \\
   & \leq |B_{n + p}| + |B_n|                             \\
   & \leq 2M
\end{align*}

综上，对$n > N$，我们有
\begin{align*}
  \left| \sum\limits_{k = n + 1}^{n + p} a_k b_k \right|
   & \leq 3\epsilon (2M) \\
   & = 6 \epsilon M
\end{align*}
故级数收敛。

\begin{zremark}
  例3：
  \begin{align*}
    2sin\frac{x}{2}\left( \frac{1}{2} + \sum\limits_{k=1}^n cos\, kx \right)
     & = sin \frac{x}{2} + \left( sin \frac{3}{2}x - sin \frac{x}{2} \right)
    + \cdots + \left[ sin(n + \frac{1}{2})x - sin(n - \frac{1}{2})x \right]
  \end{align*}
  这一步的推导。
\end{zremark}
\begin{align*}
   & 2sin\frac{x}{2}\left( \frac{1}{2} + \sum\limits_{k=1}^n cos\, kx \right)                           \\
   & = sin \frac{x}{2} + 2sin\frac{x}{2} cos x + 2sin\frac{x}{2}cos 2x + \cdots + 2sin\frac{x}{2}cos nx \\
\end{align*}
使用积化和差公式
\begin{align*}
  2 sin A cos B = sin (A + B) + sin (A - B)
\end{align*}
于是
\begin{align*}
   & sin \frac{x}{2} + 2sin\frac{x}{2} cos x + 2sin\frac{x}{2}cos 2x + \cdots + 2sin\frac{x}{2}cos nx \\                                                                      \\
   & = sin \frac{x}{2} + \left( sin \frac{3}{2}x - sin \frac{x}{2} \right)
  + \cdots + \left[ sin(n + \frac{1}{2})x - sin(n - \frac{1}{2})x \right]
\end{align*}

\begin{zremark}
  \begin{align*}
    \sum sin nx
  \end{align*}
  部分和是有界的。
\end{zremark}

习题2-(2)


\end{document}